\section{闵氏空间中的微扰匹配与欧氏空间中的非微扰重正化}

在最近一篇论文~\cite{Carlson:2017gpk} 中，把准分布因子化为部分子分布（如 Eq.~(\ref{eq:fact})）
这一做法的正确性受到了质疑。在文献~\cite{Xiong:2013bka} 原先的单圈因子化证明中，准分布与物理
部分子分布是在闵氏场论中用 Lehmann-Symanzik-Zimmermann 约化公式计算的~\cite{Sterman:1994ce}，
其中所有内部圈积分都是闵氏的。在这种形式下，所有红外（IR）性质都非常清楚。另一方面，由于准分布
没有时间依赖，它可以由欧氏空间中按
\begin{equation}
|p\rangle = e^{-\beta H}|p)
\end{equation}
演化的状态出发、在蒙特卡罗模拟中计算得到。这里 $|p\rm{)}$ 是具有动量 $p$ 和正确量子数但并非
本征态的态，$H$ 是 QCD 哈密顿量，$\beta$ 是演化参数（``欧氏时间''）。在格点 QCD 中，核子矩阵元
由源—算符—汇三点关联函数得到，其中源、算符与汇在虚时间方向上分隔开。当分隔足够大时，这个三点
关联函数由最低物理态的贡献主导~\cite{Rothe:1992nt}。因此矩阵元实际上是在一个类似欧氏的场论中
计算的。不过，欧氏计算与原闵氏场论的等价性由该形式体系的推导所保证；这一等价性历来是一切欧氏
格点计算的基础。

然而，文献~\cite{Carlson:2017gpk} 指出，带外部闵氏动量的格点微扰论费曼积分不具有共线发散。
这样一来，来自格点 QCD 的准分布（quasi-distribtutions）似乎不具备恰当的红外性质。在其后的论文
~\cite{Briceno:2017cpo} 中明确证明，格点计算流程确实给出带有正确共线发散的闵氏理论矩阵元。
问题于是变成：怎样才是从欧氏微扰论中以闵氏外动量恢复物理结果（或共线发散）的正确操作？答案是：
不能把闵氏外动量直接代入欧氏积分，因为那样会改变极点结构；正确的做法是先用欧氏外动量完成欧氏
积分，再把结果解析延拓到闵氏时空。要说明由格点 QCD 计算的准分布核子矩阵元在微扰论中确实捕捉了
正确的红外物理，就必须在欧氏圈积分完成之后再做解析延拓，而不是在此之前。

下面我们举例说明上述要点：我们计算出现在欧氏准分布中的一个单圈积分，并证明共线发散正是在从
$p_E^2$ 到 $p_M^2$ 的解析延拓之后才显现出来，其中 $p_M^\mu$ 与 $p_E^\mu$ 分别是闵氏动量与
欧氏动量。实际上结果与其闵氏对应完全一致。

我们先从闵氏积分开始，其中 $p_M^0 = \sqrt{p_z^2+m^2}$，
\beq \label{eq:mloop}
I_M = \int {d^4k_M\over (2\pi)^4}{(2\pi)p^z\delta(k^z_M-xp^z)\over (k_M^2-m^2+i\epsilon) \left[(p_M-k_M)^2+i\epsilon\right]}\ ,
\eeq
我们关注可能存在共线发散的物理区域 $0<x<1$。对于 $I_M$，先用留数定理对 $k_M^0$ 积分，再对横向
分量积分，得到
\begin{align} 	\label{eq: im}
I_M =& {i\over 8\pi}{1\over \sqrt{1+\rho_M}}\left[\ln{1+\sqrt{1+\rho_M}\over \rho_M}\right.\nonumber\\
&\left. +\ln{\sqrt{(1+\rho_M)(x^2+\rho_M)}+x+\rho_M\over 1-x}\right]\ ,
\end{align}
其中
\beq
\rho_M = {m^2\over (p^z)^2}<1\ .
\eeq

对于欧氏空间积分，先保持 $p^\tau_E$ 为实数，
\beq \label{eq:eloop}
I_E = \int {d^4k_E\over (2\pi)^4}{(2\pi)p^z\delta(k^z_E-xp^z)\over (k_E^2+m^2) (p_E-k_E)^2}  \ .
\eeq
对极点积分后我们发现
\begin{align}	\label{eq: ie}
I_E =& {1\over 8\pi\sqrt{\rho_E-1-i\ep}} \left[\arctan{\rho_E+\rho_M-2(1-x)\over2 \sqrt{\rho_E-1-i\ep}(1-x)} \right.\nonumber\\
&\left.+ \arctan {\rho_E-\rho_M-2x\over2 \sqrt{(\rho_E-1-i\ep)(x^2+\rho_M)}}\right] \ ,
\end{align}
其中
\beq \label{eq: rho}
\rho_E={(p^\tau)^2+(p^z)^2\over (p^z)^2}>1\ .
\eeq
我们在 $\rho_E$ 的虚部中保留 $-i\ep$，以便通过下半复平面、作替换 $p^\tau \rightarrow -ip^0$
把它解析延拓到 $\rho_M$，
\beq \label{eq: cont}
\sqrt{\rho_E-1-i\ep} \to -i\sqrt{1+\rho_M} \ .
\eeq
我们得到
\begin{align} \label{eq: iecont}
I'_E =&-{1\over8\pi}{1\over \sqrt{1+\rho_M}} \left[ {1\over2}\text{arctanh}\left(- {1\over\sqrt{1+\rho_M}}\right)\right.\nonumber\\
&\left. + {1\over2}\text{arctanh}\left(-{x+\rho\over\sqrt{(1+\rho_M)(x^2+\rho_M)}}\right)\right]\nonumber\\
=& -i I_M \ ,
\end{align}
这正是闵氏矩阵元，只差一个将在 $dk^0_M=idk^4_E$ 中消去的因子 $i$。若取无质量极限 $m\to0$，
则 $I_M$ 与 $I'_E$ 都给出
\beq \label{eq: oslimit}
I_M = {i\over 8\pi} \left(\ln {(p^z)^2\over m^2} + \ln {4x\over1-x} \right) \ ,
\eeq
其共线发散由 $\ln m^2$ 正确刻画。

尽管需要解析延拓才能恢复正确的共线发散，我们发现带欧氏外动量的欧氏矩阵元对重正化仍然有用。
在格点 QCD 中计算的准分布需要非微扰重正化，例如可采用文献~\cite{Alexandrou:2017huk,Chen:2017mzz,Stewart:2017tvs}
所用的 RI/MOM 方案。要在该方案中计算重正化因子，需要对具有深欧氏动量的外夸克与外胶子计算顶点
函数。重正化之后，准分布预期具有良好的连续极限，然后可在 $\overline{\rm MS}$ 方案中微扰地匹配到
通常的分布，而匹配因子可以在闵氏空间中计算。
